% !TeX root = ../../main.tex

\lstset{
    basicstyle=\ttfamily\footnotesize,
    breaklines=true,
    frame=single,
    showstringspaces=false,
    columns=fullflexible,
    literate=
        {á}{{\'a}}1
        {é}{{\'e}}1
        {í}{{\'i}}1
        {ó}{{\'o}}1
        {ú}{{\'u}}1
        {ü}{{\"u}}1
        {ñ}{{\~n}}1
}

\section{Introducción}

Este proyecto tiene como propósito configurar Apache Hadoop en Docker 
y comparar los tiempos de ejecución de un contador de palabras mediante 
MapReduce, utilizando uno y tres nodos trabajadores.

Se utilizaron los mismos programas y datos de entrada en ambas configuraciones. 
Para este proyecto, cada nodo trabajador correspondió a una pareja de servicios DataNode 
y NodeManager, además de los servicios centrales de administración del clúster.

\section{Preparación del proyecto}
Se creó una carpeta independiente llamada \texttt{2.22-proyecto-1} para organizar la 
configuración, los programas y las evidencias. 

La estructura inicial es la siguiente: 

\begin{lstlisting}[language={}]
    2.22-proyecto-1/
        docker-hadoop/
            hadoop.env
            src/
                mapper.py
                reducer.py
        evidence/
        img/
\end{lstlisting}

La carpeta \texttt{docker-hadoop} contiene los archivos del 
entorno. En \texttt{evidence} se guardaron los resultados y tiempos 
de ejecución, mientras que \texttt{img} almacenó las capturas utilizadas en el reporte. 

También se consultaron los contenedores activos mediante \texttt{docker ps}. En ese momento no había contenedores 
en ejecución.

\section{Programas de MapReduce}

Se copiaron los archivos \texttt{mapper.py} y 
\texttt{reducer.py} de la actividad anterior. Ambos programas se reutilizaron sin cambios con uno y tres trabajadores.

\noindent
\begin{minipage}{\textwidth}

\subsection{Mapper}

El mapper lee el texto desde la entrada estándar,
convierte las palabras a minúsculas y emite un par
formado por cada palabra y el valor uno.
La expresión regular permite reconocer letras,
números y caracteres como acentos y la letra ñ.

\lstinputlisting[
    language=Python,
    caption={Programa mapper para el conteo de palabras}
]{\actividadactual/docker-hadoop/src/mapper.py}

\end{minipage}
\par

Por ejemplo, para la entrada \texttt{Hadoop usa Hadoop}, 
el mapper produce:

\begin{lstlisting}[language={}]
hadoop  1
usa     1
hadoop  1
\end{lstlisting}

\noindent
\begin{minipage}{\textwidth}

\subsection{Reducer}

El reducer recibe los pares ordenados por palabra
y suma sus valores. Al cambiar de palabra, imprime
el total acumulado; al finalizar la entrada, imprime
el último grupo.

Durante la ejecución en Hadoop, el framework se
encarga de agrupar y ordenar la salida de los
mappers antes de entregarla a los reducers.

\lstinputlisting[
    language=Python,
    caption={Programa reducer para sumar las apariciones}
]{\actividadactual/docker-hadoop/src/reducer.py}

\end{minipage}
\par

Para el ejemplo anterior, el resultado esperado es:

\begin{lstlisting}[language={}]
hadoop  2
usa     1
\end{lstlisting}

\section{Configuración inicial de Hadoop}

Se creó el archivo \texttt{hadoop.env} tomando como base la configuración inicial anterior. 
Este archivo define la conexión a HDFS y los parámetros de ejecución de MapReduce mediante YARN.

Los principales valores configurados son:

\begin{center}
    \begin{tabular}{lr}
        \toprule
        \textbf{Parámetro} & \textbf{Valor} \\
        \midrule
        Factor de replicación en HDFS & 1 \\
        Memoria anunciada por trabajador & 2048 MB \\
        Memoria solicitada por tarea Map & 512 MB \\
        Memoria solicitada por tarea Reduce & 512 MB \\
        Memoria solicitada por ApplicationMaster & 512 MB \\
        \bottomrule
    \end{tabular}
\end{center}

El factor de replicación indica cuántas copias se mantienen
de cada bloque de datos; no determina el número de nodos
trabajadores.

La memoria anunciada por trabajadores es un recurso que YARN 
utiliza para planificar las tareas. Este valor no representa 
una reserva automática de memoria en Docker.

\subsection{Definición de los servicios en Docker Compose}

Se creó el archivo \texttt{docker-compose.yml} utilizando 
la imagen \texttt{ghcr.io/apache/hadoop:3.5.0} para todos los servicios.

El NameNode administra los archivos de HDFS y el 
ResourceManager coordina los recursos de ejecución. 
Los servicios DataNode almacenan los bloques de datos, 
mientras que los NodeManager permiten ejecutar las tareas.

Cada trabajador se representa mediante una pareja de 
servicios DataNode y NodeManager, ejecutados en contenedores 
separados. Todos los contenedores comparten los recursos 
de la misma computadora.

Se definieron dos configuraciones:

\begin{itemize}
    \item Un trabajador: un DataNode y un NodeManager, 
    además de los dos servicios centrales.
    En total, cuatro contenedores.
    
    \item Tres trabajadores: tres DataNode y tres NodeManager,
    además de los dos servicios centrales.
    En total, ocho contenedores.
\end{itemize}

Los trabajadores adicionales se asociaron al perfil 
\texttt{tres-nodos} de Docker Compose. Esto permite 
habilitarlos sin modificar el archivo de configuración.

Se configuraron los puertos 9870 y 8088 para acceder 
desde el navegador a las interfaces de NameNode y 
del ResourceManager, respectivamente.

\section{Prueba con un nodo trabajador}

\subsection{Inicio de los servicios}

Desde la carpeta \texttt{docker-hadoop} se inició la 
configuración de un trabajador mediante el siguiente comando:

\begin{lstlisting}[language={}]
    docker compose up -d
\end{lstlisting}

Esta configuración incluye los servicios NameNode, 
ResourceManager, DataNode y NodeManager.

Después se consultó el estado de los contenedores:

\begin{lstlisting}[language={}]
    docker compose ps -a
\end{lstlisting}

Los cuatro contenedores aparecieron con estado 
\texttt{Up}, lo que confirmó que sus procesos estaban en ejecución.

\begin{figure}[H]
    \centering
    \includegraphics[width=\textwidth]
    {\actividadactual/img/captura-1.png}
    \caption{Contenedores en ejecución para la 
    configuración de un trabajador.}
\end{figure}

\subsection{Verificación del trabajador}

Se verificó que el DataNode estuviera registrado en HDFS 
mediante el siguiente comando:

\begin{lstlisting}[language={}]
    docker compose exec -T namenode hdfs dfsadmin -report
\end{lstlisting}

El reporte mostró un DataNode activo, identificado por 
la línea \texttt{Live datanodes (1)}.

\begin{figure}[H]
    \centering
    \includegraphics[width=0.9\textwidth]
    {\actividadactual/img/captura-2.png}
    \caption{DataNode activo y registrado en HDFS.}
\end{figure}

También se verificó el registro del NodeManager en YARN:

\begin{lstlisting}[language={}]
    docker compose exec -T resourcemanager yarn node -list
\end{lstlisting}

La consulta mostró un nodo con estado \texttt{RUNNING}, 
disponible para recibir tareas.

\begin{figure}[H]
    \centering
    \includegraphics[width=\textwidth]
    {\actividadactual/img/captura-3.png}
    \caption{NodeManager registrado en YARN.}
\end{figure}

\subsection{Preparación de los datos de entrada}

Se creó el programa \texttt{generate\_data.py} para generar 
un conjunto de datos sintéticos y reproducibles. 
El texto base contiene las siguientes oraciones:

\begin{lstlisting}[language={}]
    Hadoop procesa datos.
    Docker ejecuta Hadoop.
    MapReduce cuenta palabras.
\end{lstlisting}

Este bloque se repitió 50\,000 veces en cada uno de los 
12 archivos generados. En total se obtuvieron 
1\,800\,000 líneas y 45\,000\,000 bytes, 
equivalentes a aproximadamente 42.92 MiB. 
Cada oración incluye un espacio antes del salto de línea.

Los archivos se guardaron en \texttt{src/data}. 
Se utilizó este mismo conjunto de datos en las pruebas 
con uno y tres trabajadores, sin modificar su contenido.

El uso de texto repetido permite conocer de antemano 
los conteos esperados y verificar posteriormente los 
resultados del procesamiento.

\begin{figure}[H]
    \centering
    \includegraphics[width=\textwidth]
    {\actividadactual/img/captura-4.png}
    \caption{Generación y revisión de los datos de entrada.}
\end{figure}

\subsection{Carga de los datos en HDFS}

Se creó el directorio de entrada en HDFS:
\begin{lstlisting}[language={}]
docker compose exec -T namenode \
  hdfs dfs -mkdir -p /user/hadoop/input
\end{lstlisting}

Los archivos generados estaban disponibles dentro del 
contenedor en \texttt{/workspace/data}, mediante la 
carpeta compartida configurada en Docker Compose. 
Se copiaron a HDFS con el siguiente comando:

\begin{lstlisting}[language={}]
    docker compose exec -T namenode sh -c \
        'hdfs dfs -put /workspace/data/input-*.txt /user/hadoop/input'
\end{lstlisting}

Se verificó la carga mediante los comandos 
\texttt{hdfs dfs -ls} y \texttt{hdfs dfs -count -v}. 
El directorio de entrada contenía 12 archivos, 
con un tamaño total de 45\,000\,000 bytes, 
coincidente con los datos generados localmente.

La generación y carga de los archivos forman parte 
de la preparación del entorno y no se incluyeron
en la medición del tiempo de ejecución de MapReduce.

\begin{figure}[H]
    \centering
    \includegraphics[width=\textwidth]
    {\actividadactual/img/captura-5.png}
    \caption{Archivos de entrada cargados en HDFS 
    y Verificación de su tamaño total.}
\end{figure}

\subsection{Ejecución de comprobación}

Se ejecutó el contador de palabras mediante Hadoop 
Streaming, utilizando los programas \texttt{mapper.py} 
y \texttt{reducer.py} sobre los 12 archivos de entrada.

Se configuró un reducer y se estableció que comenzara 
después de terminar todas las tareas Map. También se 
desactivaron las ejecuciones especulativas. 
Estos parámetros se mantuvieron en las pruebas 
con uno y tres trabajadores.

Esta ejecución se utilizó para comprobar el funcionamiento 
del programa y no se incluyó en los tiempos de comparación.

\begin{figure}[H]
    \centering
    \includegraphics[width=\textwidth]
    {\actividadactual/img/captura-6.png}
    \caption{Finalización correcta de la ejecución
    de comprobación con un trabajador.}
\end{figure}

El bloque de texto se repite 600\,000 veces en el conjunto 
de entrada. Por ello, el conteo esperado es de 
1\,200\,000 apariciones de la palabra \texttt{hadoop} 
y 600\,000 de cada una de las otras siete palabras.

El resultado almacenado en HDFS coincidió con
estos valores.

\begin{figure}[H]
    \centering
    \includegraphics[width=\textwidth]
    {\actividadactual/img/captura-7.png}
    \caption{Resultado del conteo de palabras 
    en la ejecución de comprobación.}
\end{figure}


\subsection{Medición del tiempo de ejecución}

Después de comprobar el resultado del programa, se
realizaron tres ejecuciones medidas con un trabajador
y se calculó el promedio de sus tiempos.

Para cada ejecución se utilizó el comando
\texttt{/usr/bin/time -p}. Se registró el valor
\texttt{real}, que representa el tiempo transcurrido
en segundos desde el inicio del comando hasta su
finalización.

La medición incluyó el envío, la espera y la ejecución
del trabajo. No incluyó la generación ni la carga
previa de los datos en HDFS.

Se mantuvieron los mismos archivos de entrada,
programas y parámetros de MapReduce. Cada ejecución
tuvo una carpeta de salida independiente y su registro
se guardó en la carpeta \texttt{evidence}.

El siguiente comando corresponde a la primera ejecución
medida de la serie utilizada para obtener el promedio
con un trabajador:

\begin{lstlisting}[language={}]
/usr/bin/time -p docker compose exec -T namenode \
  hadoop jar \
  /opt/hadoop/share/hadoop/tools/lib/hadoop-streaming-3.5.0.jar \
  -D mapreduce.job.name=wordcount-1node-series2-run1 \
  -D mapreduce.job.reduces=1 \
  -D mapreduce.job.reduce.slowstart.completedmaps=1.0 \
  -D mapreduce.map.speculative=false \
  -D mapreduce.reduce.speculative=false \
  -files /workspace/mapper.py,/workspace/reducer.py \
  -mapper "python3 mapper.py" \
  -reducer "python3 reducer.py" \
  -input /user/hadoop/input \
  -output /user/hadoop/output/1node-series2-run1 \
  2>&1 | tee ../evidence/1node-series2-run1.log
\end{lstlisting}

La opción \texttt{-files} distribuye los programas
entre los nodos que ejecutan las tareas. Las opciones
\texttt{-input} y \texttt{-output} indican las rutas
de entrada y salida en HDFS.

Se configuró un único reducer, cuyo inicio se pospuso
hasta completar todas las tareas Map. También se
desactivaron las ejecuciones especulativas.

El comando \texttt{tee} permitió mostrar los mensajes
en la terminal y guardar una copia en el archivo de
registro, incluyendo la medición del tiempo.

Para las ejecuciones restantes se cambiaron el nombre
del trabajo, la carpeta de salida y el archivo de
registro. Los parámetros de procesamiento se
mantuvieron iguales.

\subsection{Tiempos obtenidos con un trabajador}

Se realizaron tres ejecuciones medidas sobre el mismo 
conjunto de datos. Todas finalizaron correctamente. 
Los tiempos registrados corresponden al valor 
\texttt{real} del comando \texttt{/usr/bin/time -p}.


\begin{center}
    \begin{tabular}{lr}
        \toprule
        \textbf{Ejecución} & \textbf{Tiempo (s)} \\
        \midrule
        1 & 34.47 \\
        2 & 37.53 \\
        3 & 36.50 \\
        \midrule
        \textbf{Promedio} & \textbf{36.17} \\
        \bottomrule
    \end{tabular}
\end{center}

Para el promedio se utilizaron únicamente las tres 
ejecuciones de la serie \texttt{1node-series2}. 
Las ejecuciones preliminares quedaron excluidas.

\begin{figure}[H]
    \centering
    \includegraphics[width=\textwidth]
    {\actividadactual/img/captura-8.png}
    \caption{Confirmación de éxito y tiempo de las 
    tres ejecuciones con un trabajador.}
\end{figure}

\section{Prueba con tres nodos trabajadores}

\subsection{Activación de los trabajadores adicionales}

Se habilitó el perfil \texttt{tres-nodos} de Docker Compose:

\begin{lstlisting}[language={}]
    docker compose --profile tres-nodos up -d
\end{lstlisting}

Este comando inició los servicios DataNode y NodeManager 
de los trabajadores 2 y 3. Los servicios que ya estaban 
en ejecución permanecieron activos.

Se comprobó el estado de los ocho contenedores:

\begin{figure}[H]
    \centering
    \includegraphics[width=\textwidth]
    {\actividadactual/img/captura-9.png}
    \caption{Ocho contenedores activos para la 
    configuración de tres trabajadores.}
\end{figure}

\subsection{Verificación de los trabajadores}

Se consultó el estado de HDFS:

\begin{lstlisting}[language={}]
    docker compose exec -T namenode hdfs dfsadmin -report
\end{lstlisting}

El reporte mostró \texttt{Live datanodes (3)}, 
confirmando el registro de los tres DataNode.

\begin{figure}[H]
    \centering
    \includegraphics[width=\textwidth]
    {\actividadactual/img/captura-10.png}
    \caption{Tres DataNode registrados en HDFS.}
\end{figure}

Después se consultaron los nodos disponibles en YARN:

\begin{lstlisting}[language={}]
    docker compose exec -T resourcemanager yarn node -list
\end{lstlisting}

La consulta mostró tres NodeManager con estado 
\texttt{RUNNING}, disponibles para ejecutar tareas.

\begin{figure}[H]
    \centering
    \includegraphics[width=\textwidth]
    {\actividadactual/img/captura-11.png}
    \caption{Tres NodeManager registrados en YARN.}
\end{figure}

\subsection{Preparación de los datos con tres trabajadores}

Al incorporar los trabajadores adicionales, se observó 
que los nuevos DataNode todavía no contenían bloques. 
Los datos habían sido cargados cuando solo estaba disponible el primer DataNode.

Se volvieron a cargar los mismos archivos originales 
con los tres DataNode activos.

\begin{lstlisting}[language={}]
docker compose exec -T namenode sh -c \
  'hdfs dfs -put -f /workspace/data/input-*.txt /user/hadoop/input/'
\end{lstlisting}

Se mantuvieron el contenido de los archivos y el factor 
de replicación de uno.  Esta carga forma parte de la 
preparación y queda fuera de las mediciones de tiempo.

Se comprobó la distribución y el estado de los archivos 
mediante el siguiente comando:

\begin{lstlisting}[language={}]
docker compose exec -T namenode \
  hdfs fsck /user/hadoop/input -files -blocks -locations
\end{lstlisting}

\begin{center}
    \begin{tabular}{lrr}
        \toprule
        \textbf{DataNode} & \textbf{Bloques} & \textbf{Bytes} \\
        \midrule
        datanode1 & 4 & 15\,000\,000 \\
        datanode2 & 4 & 15\,000\,000 \\
        datanode3 & 4 & 15\,000\,000 \\
        \bottomrule
    \end{tabular}
\end{center}

El estado de los archivos fue \texttt{HEALTHY}, 
sin bloques perdidos ni corruptos.

\begin{figure}[H]
    \centering
    \includegraphics[width=\textwidth]
    {\actividadactual/img/captura-12.png}
    \caption{Verificación de los archivos de entrada 
    en HDFS con tres DataNode activos.}
\end{figure}

\subsection{Ejecución de comprobación}

Antes de medir los tiempos, se ejecutó el contador de 
palabras con los tres trabajadores disponibles, 
manteniendo los mismos programas y parámetros de 
MapReduce utilizados con un trabajador.

El trabajo terminó correctamente y ejecutó 12 tareas 
Map y una Reduce. El resultado coincidió con 
el obtenido en la configuración anterior: 
1\,200\,000 apariciones de \texttt{hadoop} y 
600\,000 de cada una de las otras siete palabras.

Esta ejecución previa quedó fuera del promedio 
de tiempos.


\begin{figure}[H]
    \centering
    \includegraphics[width=0.85\textwidth]
    {\actividadactual/img/captura-13.png}
    \caption{Resultado del conteo con tres
    trabajadores disponibles.}
\end{figure}

\subsection{Tiempos obtenidos con tres trabajadores}

Se realizaron tres ejecuciones medidas con los tres 
trabajadores disponibles. Todas finalizaron correctamente, 
con tiempos de 18.45, 21.41 y 21.85 segundos. 
El promedio fue de 20.57 segundos.

Se utilizó el mismo procedimiento de medición que 
con un trabajador, excluyendo la ejecución previa 
de comprobación.

Se utilizó el mismo comando de Hadoop Streaming
documentado para un trabajador, con los tres
trabajadores activos. Los nombres de las salidas
y los registros fueron \texttt{3nodes-timed1},
\texttt{3nodes-timed2} y \texttt{3nodes-timed3},
manteniendo los demás parámetros de procesamiento.

\begin{figure}[H]
    \centering
    \includegraphics[width=\textwidth]
    {\actividadactual/img/captura-14.png}
    \caption{Confirmación de éxito y tiempos de las
    tres ejecuciones con tres trabajadores.}
\end{figure}

\section{Comparación de resultados}

Las dos configuraciones procesaron los mismos 12 archivos, 
con un tamaño de 45\,000\,000 bytes. 
Se mantuvieron los programas, el factor de replicación 
y los parámetros de MapReduce, incluyendo un único reducer.

\begin{center}
    \begin{tabular}{lrr}
        \toprule
        \textbf{Ejecución} & \textbf{Un trabajador (s)} & \textbf{Tres trabajadores (s)} \\
        \midrule
        1 & 34.47 & 18.45 \\
        2 & 37.53 & 21.41 \\
        3 & 36.50 & 21.85 \\
        \midrule
        \textbf{Promedio} & \textbf{36.17} & \textbf{20.57} \\
        \bottomrule
    \end{tabular}
\end{center}

La reducción porcentual del tiempo se calculó mediante:

\[
    \text{reducción} = 
    \frac{\overline{T}_1 - \overline{T}_3}
    {\overline{T}_1}
    \times 100 
    \approx 43.12\text{\,\%}
\]

Donde $\overline{T}_1$ y $\overline{T}_3$ representan 
los promedios con uno y tres trabajadores, respectivamente. 
El porcentaje se calculó usando los promedios sin redondear.

La configuración de tres trabajadores redujo el tiempo 
aproximadamente 15.60 segundos.

MapReduce permite distribuir las tareas Map entre los trabajadores disponibles. 
Al incorporar más trabajadores, YARN dispone de más 
recursos para ejecutar tareas simultáneamente. 
Sin embargo, en esta prueba se mantuvo un único reducer y 
la medición también incluyó el envío y la coordinación 
del trabajo.

Todos los contenedores compartieron los recursos de 
la misma computadora. Además, se utilizó un conjunto 
de datos sintéticos de aproximadamente 42.92 MiB. 
Por ello, los resultados corresponden a las condiciones 
de esta práctica y no representan una mejora garantizada 
para cualquier volumen de datos o equipo.

\section{Conclusión}

Se configuró Apache Hadoop en Docker con uno y tres 
nodos trabajadores y se ejecutó un contador de palabras 
en Python mediante Hadoop Streaming.

Las ejecuciones de comprobación produjeron el mismo 
conteo en ambas configuraciones. En las pruebas medidas, 
el tiempo promedio disminuyó de 36.17 a 20.57 segundos, 
lo que representa una reducción del 43.12\,\%.

La práctica permitió comprobar el funcionamiento de HDFS 
y YARN, verificar la distribución de los datos y comparar 
el tiempo de ejecución de MapReduce al aumentar 
la cantidad de trabajadores.